\documentclass[a4paper,12pt,titlepage,finall]{article}

\usepackage[T1,T2A]{fontenc}     % форматы шрифтов
\usepackage[utf8x]{inputenc}     % кодировка символов, используемая в данном файле
\usepackage[russian]{babel}      % пакет русификации
\usepackage{tikz}                % для создания иллюстраций
\usepackage{pgfplots}            % для вывода графиков функций
\usepackage{geometry}		 % для настройки размера полей
\usepackage{indentfirst}         % для отступа в первом абзаце секции
\usepackage{multirow}            % для таблицы с результатами

% выбираем размер листа А4, все поля ставим по 3см
\geometry{a4paper,left=30mm,top=30mm,bottom=30mm,right=30mm}

\setcounter{secnumdepth}{0}      % отключаем нумерацию секций

\usepgfplotslibrary{fillbetween} % для изображения областей на графиках

\begin{document}
% Титульный лист
\begin{titlepage}
    \begin{center}
	{\small \sc Московский государственный университет \\имени М.~В.~Ломоносова\\
	Факультет вычислительной математики и кибернетики\\}
	\vfill
	{\Large \sc Отчет по заданию №1}\\
	~\\
	{\large \bf <<Методы сортировки>>}\\ 
	~\\
	{\large \bf Вариант 5 / 5/ 5/ 5}
    \end{center}
    \begin{flushright}
	\vfill {Выполнил:\\
	студент 104 группы\\
	Ильин А.Д.\\
	~\\
	Преподаватель:\\
	Гуляев Д.А.}
    \end{flushright}
    \begin{center}
	\vfill
	{\small Москва\\2026}
    \end{center}
\end{titlepage}

% Автоматически генерируем оглавление на отдельной странице
\tableofcontents
\newpage

\section{Постановка задачи}

В данной работе требуется реализовать два алгоритма сортировки одномерного массива целых чисел — пузырьковую сортировку и быструю сортировку с разбиением Ломуто — и провести экспериментальное сравнение их эффективности. Оба метода должны упорядочивать элементы по возрастанию.

Для исследования необходимо создать несколько вариантов исходных данных:
\begin{itemize}
\item полностью упорядоченный массив;
\item массив, отсортированный в обратном порядке;
\item два различных массива со случайным расположением элементов (заполняются с помощью генератора случайных чисел).
\end{itemize}

Сравнение алгоритмов выполняется на одних и тех же наборах данных при разных размерах массива: 10, 100, 1000 и 10000 элементов. Для каждого размера и типа исходного массива проводится несколько запусков (например, 5), после чего усредняются две ключевые метрики:
\begin{itemize}
\item количество сравнений элементов;
\item количество перемещений (обменов или присваиваний) элементов в процессе сортировки.
\end{itemize}

Полученные результаты позволяют оценить зависимость трудоёмкости алгоритмов от объёма и характера входных данных и сделать выводы об их применимости в различных условиях.

\newpage

\section{Результаты экспериментов}

В ходе эксперимента были получены следующие данные.

Пузырьковая сортировка. 
Количество сравнений для всех типов массивов при фиксированном \(n\) остаётся неизменным и равно \( \frac{n(n-1)}{2} \), что полностью соответствует теоретической оценке. Количество перемещений зависит от исходной упорядоченности:
\begin{itemize}
\item на упорядоченном массиве перемещения отсутствуют;
\item на полностью обратном массиве число перемещений максимально и также равно \( \frac{n(n-1)}{2} \);
\item на случайных массивах число перемещений принимает промежуточные значения (например, при \(n = 10\) в среднем около 20–25).

Быстрая сортировка с разбиением Ломуто. 
Для случайных массивов число сравнений и перемещений растёт пропорционально \(O(n \log n)\), что заметно при увеличении \(n\): так, при переходе от \(n = 1000\) к \(n = 10000\) рост составляет примерно 10–12 раз, а не 100 раз, как в случае квадратичной зависимости. На упорядоченных массивах (возрастающем и убывающем) наблюдается худший случай: количество сравнений приближается к \( \frac{n(n-1)}{2} \), что соответствует квадратичной сложности. Это объясняется выбором последнего элемента в качестве опорного.

Полученные экспериментальные данные подтверждают теоретические оценки рассматриваемых алгоритмов сортировки.

\begin{table}[h]
\centering
\begin{tabular}{|c|c|c|c|c|c|c|c|}
    \hline
    \multirow{2}{*}{\textbf{n}} & \multirow{2}{*}{\textbf{Параметр}} & \multicolumn{4}{|c|}{\textbf{Номер сгенерированного массива}} & \textbf{Среднее} \\
    \cline{3-6}
    & & \parbox{1.5cm}{\centering 1} & \parbox{1.5cm}{\centering 2} & \parbox{1.5cm}{\centering 3} & \parbox{1.5cm}{\centering 4} & \textbf{значение} \\
    \hline
    \multirow{2}{*}{10} & Сравнения & 45& 45& 27& 23& 35\\
    \cline{2-7}
                        & Перемещения & 0& 5& 11& 6& 5.5\\
    \hline
    \multirow{2}{*}{100} & Сравнения & 4950& 4950& 606& 672& 2794.5\\
    \cline{2-7}
                         & Перемещения & 0& 50& 324& 272& 161.5\\
    \hline
    \multirow{2}{*}{1000} & Сравнения & 499500& 499500& 10550& 11208& 255189.5\\
    \cline{2-7}
                          & Перемещения & 0& 500& 5038& 5588& 2781.5\\
    \hline
    \multirow{2}{*}{10000} & Сравнения & 49995000& 49995000& 174260& 179592& 25085963\\
    \cline{2-7}
                           & Перемещения & 0& 5000& 67004& 71881& 35971.25\\
    \hline
\end{tabular}
\caption{Результаты работы сортировки Quick sort (Lomuto)}
\end{table}

\newpage

\section{Структура программы и спецификация функций}
\begin{itemize}
\item genOrdered(int a[], int n) – заполняет массив a упорядоченными по возрастанию числами от 1 до n.
\item genReversed(int a[], int n) – заполняет массив a числами в обратном порядке: n, n‑1, …, 1.
\item genRand1(int a[], int n) – заполняет массив a случайными числами от 0 до 999 (первый вариант случайного массива).
\item genRand2(int a[], int n) – заполняет массив a случайными числами от 0 до 999 (второй вариант случайного массива).
\item bubbleSort(int a[], int n) – реализует пузырьковую сортировку массива a по возрастанию; в процессе работы увеличивает глобальные счётчики `cmp` (число сравнений) и `mov` (число обменов).
\item partition(int a[], int lo, int hi) – статическая вспомогательная функция для быстрой сортировки; выполняет разбиение подмассива a[lo..hi] по схеме Ломуто (опорный элемент – последний). Возвращает индекс, на который опорный элемент становится после разбиения. Увеличивает `cmp` при каждом сравнении и `mov` при каждом обмене.
\item qs(int a[], int lo, int hi) – статическая рекурсивная функция быстрой сортировки; сортирует подмассив a[lo..hi] с помощью вызова `partition` и рекурсивных вызовов для левой и правой частей.
\item quickSort(int a[], int n) – обёртка для вызова быстрой сортировки всего массива; вызывает `qs` с границами 0 и n‑1.
\item copyArray(int dest[], int src[], int n) – копирует содержимое массива `src` в массив `dest`; используется для создания копий исходных массивов перед каждым запуском сортировки.

\newpage

\section{Отладка программы, тестирование функций}

Отладка и тестирование алгоритмов проводились с помощью автоматической проверки упорядоченности массива после сортировки. Для этого использовалась вспомогательная функция:
\begin{verbatim}
int is_sorted(const int *arr, int n) {
    for (int i = 1; i < n; ++i) {
        if (arr[i] < arr[i - 1]) {
            return 0;
            }
    }
    return 1;
}
\end{verbatim}

Рассматривались следующие случаи:
\begin{itemize}
\item Граничные случаи: пустой массив (n = 0) и массив из одного элемента. Оба алгоритма корректно обработали эти ситуации (сортировка не вносила изменений, функция `is_sorted` возвращала 1).
\item Уже отсортированный массив по возрастанию: например, {1, 2, 3, 4, 5}. После сортировки массив остался неизменным, `is_sorted = 1`.
\item Массив, отсортированный по убыванию (обратный порядку): {5, 4, 3, 2, 1}. Оба метода успешно переупорядочили его в {1, 2, 3, 4, 5}.
\item Случайная последовательность: например, {3, 1, 4, 1, 5, 9, 2, 6, 5}. После сортировки массив удовлетворял условию возрастания.
\end{itemize}
Для каждого теста выполнялось по несколько запусков, в том числе на больших массивах (до 10000 элементов). Все результаты подтвердили корректность реализованных алгоритмов.

\newpage

\section{Анализ допущенных ошибок}

В процессе разработки были выявлены и исправлены следующие ошибки:
\begin{enumerate}
\item Неправильный подсчёт числа перемещений. Первоначально при каждом обмене двух элементов счётчик `mov` увеличивался на 3 (по числу присваиваний). Однако в соответствии с заданием перемещением считается один обмен, а не отдельные присваивания. После исправления операция обмена стала увеличивать счётчик только на 1.

\item Лишние обмены в быстрой сортировке при разбиении. В реализации схемы Ломуто при `i != j` обмен не требовался, но в исходном коде обмен выполнялся всегда, что приводило к завышению числа перемещений на уже упорядоченных массивах. Добавлена проверка `if (i != j)` перед обменом, а также аналогичная проверка при установке опорного элемента на место, чтобы избежать обмена элемента с самим собой.

\itemСброс глобальных счётчиков между запусками. Изначально забывали обнулять `cmp` и `mov` перед каждым вызовом сортировки, что приводило к накоплению статистики от предыдущих запусков. Ошибка устранена обнулением счётчиков непосредственно перед вызовом сортирующей функции.
\end{enumerate}
Все перечисленные ошибки были исправлены, итоговая программа работает корректно во всех предусмотренных режимах.

\newpage
\begin{raggedright}
\addcontentsline{toc}{section}{Список цитируемой литературы}
\begin{thebibliography}{99}
\item Кормен Т., Лейзерсон Ч., Ривест Р., Штайн К. Алгоритмы: построение и анализ. Второе издание. — М.: «Вильямс», 2005.
\item Вирт Н. Алгоритмы и структуры данных. — М.: Мир, 1989.
\item Трифонов Н. П., Пильщиков В. Н. Задания практикума на ЭВМ (1 курс). Методическая разработка. — М.: ВМК МГУ, 2001.
\end{thebibliography}
\end{raggedright}

\end{document}
